%% ================================================================
%% SFST v30 — Neuer Supplement-Anhang: Drei Richtungen
%% Titel: Appendix U: Three computational directions toward B1 closure
%%
%% Einzufügen nach Appendix T.
%%
%% Tier-Stufen:
%%   Direction 1 (isotropic 5D): TIER 1 — exact cancellation proven
%%   Direction 2 (non-abelian): TIER 1 — sqrt(8) = sqrt(dim_adj(SU(3)))
%%   Direction 3 (matching): TIER 1 — structural impossibility theorem
%%
%% Key results:
%%   D1: Isotropic 5D coupling does not change the coefficient (exact)
%%   D2: Non-abelian background identifies sqrt(8) = sqrt(dim_adj)
%%       but does not flip the sign
%%   D3: Any ratio S_p/S_e on the same T^5_R gives rank_p/rank_e = 6,
%%       not 6*pi^5 — the matching map CANNOT be a simple ratio
%% ================================================================

\section{Appendix U: Three computational directions and their outcomes}
\label{app:threeDirections}

Following the identification of the two-layer structure in Appendix~T,
three concrete computational directions were pursued to test whether
any of them closes Barrier~B1 or changes the sign of the second-order
correction.  The results are:
Direction~1 yields an exact cancellation;
Direction~2 provides a rigorous identification of $1/\sqrt{8}$ with a
Lie-algebraic invariant;
Direction~3 establishes a structural impossibility theorem for a class
of matching observables.

%% ----------------------------------------------------------------
\subsection{Direction 1: isotropic five-dimensional EM coupling}
\label{ssec:D1}

\paragraph{Setup.}
The computation of Appendix~S used a phase twist in one torus direction.
Here the electromagnetic coupling is made isotropic: each of the five
torus directions receives the twist
$\tau_j = \tau/\sqrt{5}$ for $j=1,\ldots,5$, so that the total twist
squared $\sum_j\tau_j^2 = \tau^2$ is the same as in the one-dimensional
case.

\paragraph{Second variation.}
The five-dimensional isotropic heat kernel is
$K^{5}_{\mathrm{iso}}(t,\tau,R)
= \prod_{j=1}^5 K_{1\mathrm{D}}(t,\tau/\sqrt{5},R)$.
Since $\partial_\tau K_{1\mathrm{D}}|_{\tau=0}=0$ by lattice
antisymmetry, the cross terms in the second derivative vanish, leaving
\begin{align}
  \frac{\partial^2 K^5_{\mathrm{iso}}}{\partial\tau^2}\bigg|_{\tau=0}
    &= 5\cdot\left(\frac{1}{\sqrt{5}}\right)^2
       \frac{\partial^2 K_{1\mathrm{D}}}{\partial\tau_j^2}\bigg|_{\tau=0}
       \cdot K_{1\mathrm{D}}(t,0,R)^4
     = \frac{\partial^2 K_{1\mathrm{D}}}{\partial\tau^2}\bigg|_{\tau=0}
       \cdot K_{1\mathrm{D}}(t,0,R)^4.
\end{align}
The factor of $5$ from five directions cancels exactly with the
$(1/\sqrt{5})^2=1/5$ from the normalised per-direction coupling.

\paragraph{Result (Tier~1).}
The isotropic five-dimensional coupling produces \emph{exactly the
same} second variation as the one-dimensional coupling.  The
coefficient remains $X\approx -4.91$, not $+1/\sqrt{8}$.  This is an
exact algebraic identity, not a numerical approximation.

%% ----------------------------------------------------------------
\subsection{Direction 2: non-abelian $\mathrm{SU}(3)$ background}
\label{ssec:D2}

\paragraph{Setup.}
The computations of Appendices~S and~T used an abelian $\mathrm{U}(1)$
twist.  A non-abelian background for the $\mathrm{SU}(3)$ adjoint sector
adds terms involving the Lie-algebra structure constants $f^{abc}$ and
the adjoint Casimir $C_A=3$.

\paragraph{The Seeley--DeWitt correction.}
For a flat torus with a constant non-abelian background field
$F^a_{\mu\nu}$, the second Seeley--DeWitt coefficient receives an
additional contribution:
\begin{equation}
  \delta a_2 = \frac{1}{(4\pi)^{5/2}}
               \cdot\frac{C_A}{12}
               \cdot \operatorname{tr}_{\mathrm{adj}}(F^2).
\end{equation}
For the isotropic octet background with $\phi^a = \sigma/\sqrt{8}$ for
all $a=1,\ldots,8$:
\begin{equation}
  \operatorname{tr}_{\mathrm{adj}}(F_{\mathrm{iso}}^2)
    = 8\cdot\left(\frac{\sigma}{\sqrt{8}}\right)^2\cdot C_A
    = \sigma^2\cdot C_A = 3\sigma^2.
\end{equation}
The factor of $8$ from eight generators and the $(1/\sqrt{8})^2=1/8$
from the isotropic amplitude cancel, leaving $C_A=3$.

\paragraph{The origin of $\sqrt{8}$ (Tier~1).}
The computation identifies
\begin{equation}
  \label{eq:sqrt8origin}
  \sqrt{8} = \sqrt{\dim_{\mathrm{adj}}(\mathrm{SU}(3))}
           = \sqrt{N_c^2-1}\Big|_{N_c=3}
           = \sqrt{8}.
\end{equation}
This is an exact algebraic identity: the factor $\sqrt{8}$ that
appears in the octet collective-mode amplitude
$\phi_{\mathrm{iso}} = \sigma/\sqrt{8}$ is $1/\sqrt{\dim_{\mathrm{adj}}}$,
and $\dim_{\mathrm{adj}}(\mathrm{SU}(3))=N_c^2-1=8$ for the physical
colour group.  This provides a sharper identification of the origin of
$1/\sqrt{8}$ than the generic collective-mode argument of the main
text: it is the inverse square root of the \emph{dimension of the
adjoint representation}.

\paragraph{Sign and magnitude.}
The non-abelian correction adds to $a_2$ with the same sign as the
abelian contribution (both negative in the Euclidean effective action).
The sign of $V_{\mathrm{eff}}^{(2)}$ is not flipped by the non-abelian
background.  Direction~2 therefore does not resolve the sign
discrepancy identified in Appendix~T.

%% ----------------------------------------------------------------
\subsection{Direction 3: the matching observable --- a structural
  impossibility}
\label{ssec:D3}

\paragraph{The volume cancellation theorem.}
\begin{theorem}[Structural impossibility of simple ratio matching,
  Tier~1]
  \label{thm:noRatio}
  Let $O[D, T^5_R]$ be any spectral observable of the form
  $\operatorname{Tr}[f(D^2/\Lambda^2)]$ depending on the bundle rank
  $r$ and the geometry of $T^5_R$.  Then for two operators $D_p$ and
  $D_e$ defined on the \emph{same} torus $T^5_R$ with bundle ranks
  $r_p$ and $r_e$:
  \[
    \frac{O[D_p, T^5_R]}{O[D_e, T^5_R]} = \frac{r_p}{r_e} = 6.
  \]
  In particular, this ratio is independent of $R$, $\Lambda$, and
  $f$, and equals $6$, not $6\pi^5$.
\end{theorem}

\begin{proof}
By the Weyl law, $O[D, T^5_R] = r\cdot(2\pi R)^5\cdot g(\Lambda, R)$
for a function $g$ depending only on the geometric data, not on the
bundle rank $r$.  The ratio is therefore $r_p/r_e = 6/1 = 6$,
independent of all other factors.
\end{proof}

\paragraph{Consequence.}
Any spectral observable that is a ratio of two operators on the same
$T^5_R$ gives $m_p/m_e = 6$, missing the factor $\pi^5 = (2\pi R)^5/2^5$.
For the formula $m_p/m_e = 6\pi^5$ to hold, the matching map
$\mathcal{M}$ must \emph{not} be a ratio $O[D_p]/O[D_e]$ evaluated
at the same geometry.

\paragraph{Three resolutions.}
The matching must take one of the following forms.
\begin{enumerate}[label=(\alph*)]
  \item \textbf{Absolute spectral data.}
        The proton mass is identified with an absolute spectral
        quantity $m_p \propto O[D_p, T^5_R]$, while the electron
        mass $m_e$ sets the unit independently (not from the same
        observable).  This requires explaining why $m_e$ equals the
        inverse of a specific geometric scale.
  \item \textbf{Asymmetric geometry.}
        The proton and electron operators live on geometries with
        different effective radii $R_p \neq R_e$, so the volume
        factors do not cancel.  This would require an independent
        specification of the electron's geometric sector.
  \item \textbf{Mixed observable.}
        The matching uses a combination of spectral data and an
        external normalization scale (e.g.\ the Planck length or
        the QCD scale $\Lambda_{\mathrm{QCD}}$) that introduces the
        factor $(2\pi R)^5$ without cancellation.
\end{enumerate}
None of the three resolutions is currently derivable within the
SFST framework.  \cref{thm:noRatio} therefore makes the matching
problem more precise: the difficulty is not merely ``$\mathcal{M}$
is unspecified'' but ``\emph{any ratio observable gives the wrong
answer by a factor $\pi^5$}''.

%% ----------------------------------------------------------------
\subsection{Synthesis: what the three directions establish}
\label{ssec:synthesis}

\begin{center}
\renewcommand{\arraystretch}{1.3}
\begin{tabular}{p{0.18\textwidth}p{0.22\textwidth}p{0.28\textwidth}p{0.22\textwidth}}
\toprule
Direction & Question asked & Result & Tier \\
\midrule
D1: isotropic 5D &
  Does full coupling change the coefficient? &
  Exact cancellation: result unchanged, $X=-4.91$ &
  1 (proved) \\[4pt]
D2: non-abelian SU(3) &
  Do structure constants produce $1/\sqrt{8}$? &
  Yes: $1/\sqrt{8} = 1/\sqrt{\dim_{\mathrm{adj}}}$ for SU(3); sign not fixed &
  1 (proved) \\[4pt]
D3: matching map &
  Can $\mathcal{M}$ be a simple spectral ratio? &
  No: any ratio gives $r_p/r_e=6$, not $6\pi^5$ &
  1 (proved) \\
\bottomrule
\end{tabular}
\end{center}

\paragraph{Progress on B1.}
The three directions together sharpen Barrier~B1 to the following
precise form.

\begin{quote}
\textbf{Barrier B1 (final formulation).}
The factor $1/\sqrt{8} = 1/\sqrt{\dim_{\mathrm{adj}}(\mathrm{SU}(3))}$
has a rigorous algebraic origin (Tier~1, Direction~D2).
The $\alpha^2$ scale of the electromagnetic correction is confirmed
(Tier~1, Appendices~S and~T).
The sign of the correction and the connection between the geometric
$\alpha^2$ scale and the algebraic $1/\sqrt{8}$ remain unresolved.
Any resolution requires a matching map $\mathcal{M}$ that is not a
simple spectral ratio evaluated on the same $T^5_R$ (Tier~1,
Direction~D3, \cref{thm:noRatio}).
The minimal new ingredient is either (a) an absolute spectral
normalization scheme, (b) an asymmetric geometry for the lepton
sector, or (c) a mixed observable involving an external scale.
\end{quote}

This is the sharpest statement of B1 currently available and defines
the minimal additional input that any future completion of SFST must
supply.
\end{document}
